\chapter{Markovian Queues: M/M/1, Processor Sharing, and Insensitivity}\label{sec:L2}

\Cref{sec:L1} dealt with facts, such as the arrival process and Little's law, that
depend hardly at all on the details of the model. This lecture treats the case where
the state of the system (the number of jobs present) is a continuous-time Markov
chain. The stationary distribution can then be computed exactly, and the model of the
paper's decode stage (processor sharing) and its key property, \emph{insensitivity},
come out of it.

\section{Continuous-time Markov chains}

\begin{definition}[Continuous-time Markov chain]\label{def:ctmc}
A \term{generator} ($Q$-matrix) on a countable set $\mathcal{X}$ is a matrix
$Q=(q(x,y))_{x,y\in\mathcal{X}}$ with $q(x,y)\ge 0$ for $x\ne y$,
$q(x):=\sum_{y\ne x}q(x,y)<\infty$, and $q(x,x)=-q(x)$.
A right-continuous process $(X_t)_{t\ge0}$ with values in $\mathcal{X}$ is a
\term{continuous-time Markov chain} (CTMC) with generator $Q$ if, on entering state
$x$, it stays there for an exponential time with mean $1/q(x)$ (independent of the
past) and then jumps to $y$ with probability $q(x,y)/q(x)$. $q(x,y)$ is the
\term{transition rate} from $x$ to $y$.
The chain is \term{non-explosive} if infinitely many jumps in finite time have
probability zero, and \term{irreducible} if every $y$ can be reached from every $x$
with positive probability.
\end{definition}

The memorylessness of the exponential distribution (\cref{prop:memoryless}) is what
makes this definition possible. If holding times were not exponential, how long the
chain has already stayed would affect the future and the Markov property would fail.
Intuitively, $q(x,y)$ means that the probability of moving from $x$ to $y$ in a short
time $h$ is $q(x,y)h+o(h)$.

\begin{definition}[Stationary distribution, balance equations]\label{L2:def:balance}
A probability distribution $\pi$ on $\mathcal{X}$ with $\pi Q=0$, that is, for all $x$,
\begin{equation}\label{L2:eq:global}
  \pi(x)\,q(x)=\sum_{y\ne x}\pi(y)\,q(y,x),
\end{equation}
is a \term{stationary distribution}, and \eqref{L2:eq:global} are the
\term{global balance equations}. The left side is the probability flow out of $x$,
the right side the probability flow into $x$. The stronger condition
\begin{equation}\label{L2:eq:detailed}
  \pi(x)\,q(x,y)=\pi(y)\,q(y,x)\qquad(\forall x\ne y)
\end{equation}
is called \term{detailed balance}.
\end{definition}

\begin{lemma}[Detailed balance \cite{norris1997}]\label{L2:lem:detailed}
A probability distribution satisfying detailed balance is stationary.
\end{lemma}
\begin{proof}
Summing \eqref{L2:eq:detailed} over $y\ne x$ gives $\pi(x)\sum_{y\ne x}q(x,y)=\sum_{y\ne x}\pi(y)q(y,x)$, which is \eqref{L2:eq:global}.
\end{proof}

We quote the next theorem without proof \cite[Ch.~3]{norris1997}. It lets us read the
stationary probability $\pi(x)$ as ``the fraction of time the state is $x$''.

\begin{theorem}[Ergodic theorem \cite{norris1997}]\label{L2:thm:ergodic}
If an irreducible, non-explosive CTMC has a stationary distribution $\pi$, then $\pi$
is unique, the chain is positive recurrent, and for any initial state and any $f$ with
$\sum_x|f(x)|\pi(x)<\infty$, almost surely
\[
  \frac1t\int_0^t f(X_s)\dd s\ \longrightarrow\ \sum_x f(x)\pi(x)\qquad(t\to\infty).
\]
\end{theorem}

\section{Birth--death processes and M/M/1}

\begin{definition}[Birth--death process]\label{L2:def:bd}
A CTMC on $\mathcal{X}=\{0,1,2,\dots\}$ whose only nonzero transition rates are
$q(n,n+1)=\lambda_n>0$ ($n\ge0$) and $q(n,n-1)=\mu_n>0$ ($n\ge1$) is a
\term{birth--death process}. $\lambda_n$ are the birth rates and $\mu_n$ the death
rates.
\end{definition}

In a queue, $n$ is the number of jobs in the system. An arrival increases $n$ by one
and a completion decreases it by one, so if arrivals and completions happen on
exponential clocks, the number of jobs is a birth--death process. It is known that
the process is non-explosive when the birth rates are bounded.

\begin{theorem}[Stationary distribution of a birth--death process \cite{kleinrock1975}]\label{thm:birthdeath}
A birth--death process has a stationary distribution if and only if
\begin{equation}\label{L2:eq:bdsum}
  H:=\sum_{n=0}^{\infty}\prod_{k=1}^{n}\frac{\lambda_{k-1}}{\mu_k}<\infty
\end{equation}
(an empty product is 1), and then the stationary distribution is
$\pi(n)=H^{-1}\prod_{k=1}^{n}\lambda_{k-1}/\mu_k$.
\end{theorem}
\begin{proof}
(Necessity and the form of $\pi$.) Suppose $\pi$ satisfies \eqref{L2:eq:global}. At
state $0$, $\pi(0)\lambda_0=\pi(1)\mu_1$. For $n\ge1$, assume
$\pi(n-1)\lambda_{n-1}=\pi(n)\mu_n$. Cancelling this term from the global balance at
state $n$, $\pi(n)(\lambda_n+\mu_n)=\pi(n-1)\lambda_{n-1}+\pi(n+1)\mu_{n+1}$, gives
$\pi(n)\lambda_n=\pi(n+1)\mu_{n+1}$. By induction, for all $n$
\begin{equation}\label{L2:eq:cut}
  \pi(n)\lambda_n=\pi(n+1)\mu_{n+1},
\end{equation}
hence $\pi(n)=\pi(0)\prod_{k=1}^n\lambda_{k-1}/\mu_k$. For $\sum_n\pi(n)=1$ we need
$\pi(0)=1/H>0$, that is, $H<\infty$.
(Sufficiency.) If $H<\infty$, the $\pi$ defined above is a probability distribution,
and \eqref{L2:eq:cut} is exactly detailed balance \eqref{L2:eq:detailed}, so $\pi$ is
stationary by \cref{L2:lem:detailed}.
\end{proof}

\Cref{L2:eq:cut} is called the \term{cut equation}: the flow across the boundary
between $\{0,\dots,n\}$ and $\{n+1,\dots\}$ is the same in both directions.

\begin{definition}[M/M/1 queue]\label{L2:def:mm1}
In Kendall notation, \term{M/M/1} is the queue with Poisson arrivals of rate $\lambda$
(\cref{def:poisson}), service times exponential with mean $1/\mu$ (independent of each
other and of the arrivals), one server, and unlimited waiting room. Its load is
$\rho=\lambda/\mu$ (\cref{def:load}).
\end{definition}

\begin{proposition}[M/M/1 \cite{kleinrock1975}]\label{prop:mm1}
In M/M/1, the number of jobs $N_t$ (including the one in service) is a birth--death
process with $\lambda_n=\lambda$, $\mu_n=\mu$. If $\rho<1$, the stationary
distribution is geometric, $\pi(n)=(1-\rho)\rho^n$, and
\[
  L=\frac{\rho}{1-\rho},\qquad W=\frac{1}{\mu-\lambda},\qquad W_q=\frac{\rho}{\mu(1-\rho)} .
\]
If $\rho\ge1$, there is no stationary distribution. As $\rho\uparrow1$, $L$, $W$ and
$W_q$ all diverge to $+\infty$.
\end{proposition}
\begin{proof}
When $N_t=n\ge1$, by memorylessness the time to the next arrival is $\mathrm{Exp}(\lambda)$,
the residual service time of the job in service is $\mathrm{Exp}(\mu)$, and the two are
independent. So the rate of $n\to n+1$ is $\lambda$ and that of $n\to n-1$ is $\mu$.
In \cref{thm:birthdeath}, $H=\sum_n\rho^n$ is finite only when $\rho<1$, and then
$H=1/(1-\rho)$. Hence $\pi(n)=(1-\rho)\rho^n$ and
$L=\sum_n n(1-\rho)\rho^n=\rho/(1-\rho)$. By Little's law (\cref{thm:little}),
$W=L/\lambda=1/(\mu-\lambda)$, and since the waiting time is the sojourn time minus the
service time, $W_q=W-1/\mu=\rho/(\mu(1-\rho))$.
\end{proof}

At $\rho=0.9$ we get $L=9$, and at $\rho=0.99$, $L=99$. This factor $1/(1-\rho)$, which
says that near capacity a small increase in mean load greatly increases the delay,
reappears in the ``price of a miss'' of \cref{sec:L3}.

\section{Processor sharing and state-dependent capacity}

So far the server has processed one job at a time. The decode stage of an LLM does not.

\begin{definition}[Work, processor sharing]\label{def:ps}
The \term{work} (service requirement) $S\ge0$ of a job is the time the server needs to
process it at full speed.
Given a \term{capacity function} $\varphi:\{1,2,\dots\}\to(0,\infty)$, the rule under
which, when $n$ jobs are present, the server processes $\varphi(n)$ units of work per
second and splits it equally among the $n$ jobs, so that each job's remaining work
decreases at rate $\varphi(n)/n$, is called \term{processor sharing} (PS).
$\varphi\equiv1$ is classical PS; with Poisson arrivals and a general work distribution
$G$ we write \term{M/G/1-PS}. The load is $\rho=\lambda\E[S]$.
\end{definition}

\serving{In continuous batching, one decode step advances every request in the batch
by one token. We therefore model the decode stage as PS, with the $m$ requests sharing
the device equally. Because the weights are read once for the whole batch, $\varphi(m)$
increases with $m$, and because memory (the KV cache) limits the batch size, it
eventually saturates. The decode stage of \S2.2 of the paper is exactly this model:
the work is the decode demand $D_i$ and the load is $\rho_D=\lambda\E[D]$.}

\begin{proposition}[PS with exponential work \cite{kelly1979}]\label{prop:phi}
For a PS server with capacity $\varphi$, Poisson arrivals of rate $\lambda$, and
exponential work with mean $\E[S]$, the number of jobs is a birth--death process with
$\lambda_n=\lambda$ and $\mu_n=\varphi(n)/\E[S]$. Hence, with
$F(n):=\prod_{k=1}^n\varphi(k)$, if $G(\rho):=\sum_{n\ge0}\rho^n/F(n)<\infty$ then
\begin{equation}\label{L2:eq:piphi}
  \pi(n)=\frac{\rho^n/F(n)}{G(\rho)},\qquad \rho=\lambda\E[S].
\end{equation}
\end{proposition}
\begin{proof}
In state $n$, by memorylessness the remaining works of the $n$ jobs are independent,
each $\mathrm{Exp}(1/\E[S])$, and each decreases at speed $\varphi(n)/n$.
Hence each job's completion time is exponential with rate $\varphi(n)/(n\E[S])$, and the
first of the $n$ completions is exponential with rate $\varphi(n)/\E[S]$
(\cref{L1:lem:min}). The rest follows from \cref{thm:birthdeath}, since
$\prod_{k\le n}\lambda\E[S]/\varphi(k)=\rho^n/F(n)$.
\end{proof}

With $\varphi\equiv1$, \eqref{L2:eq:piphi} is the geometric distribution of M/M/1. With
exponential work, the service order does not change the distribution of the number of
jobs.

\begin{example}[Saturating capacity]\label{L2:ex:sat}
Let $\varphi(m)=\min(m,c)$. Below $c$ the throughput grows in proportion to the batch
size, and it saturates at $c$ (the simplest model of a memory that holds only $c$
requests). This is the same birth--death process as M/M/$c$ with $c$ servers, and by
\eqref{L2:eq:piphi}, $\pi(n)\propto\rho^n/n!$ for $n\le c$ and
$\pi(n)\propto\rho^n/(c!\,c^{n-c})$ for $n\ge c$.
$G(\rho)<\infty$ if and only if $\rho<c=\sup\varphi$. If the load exceeds the saturated
capacity, the system is unstable however large the batch grows.
\end{example}

The next theorem says that more expected demand never leaves fewer jobs. This sounds
obvious, but it is not true in general: it depends on the discipline.

\paragraph{More work, fewer jobs.} At a single FIFO server compare two work laws. Under
law A a job takes $0.1$ with probability $0.9$ and $9.1$ with probability $0.1$; under
law B every job takes $1.2$, so B carries $20\,\%$ more work. In \cref{fig:more-work}
the long job of law A holds the server while short jobs pile up behind it, and the
number in system reaches $7$; under law B it never exceeds $2$. In steady state, with
Poisson arrivals at $\lambda=0.5$ ($\rho=0.5$ for A, $0.6$ for B), the
Pollaczek--Khinchine formula of \cref{sec:L3} gives a mean number in system of $2.57$
under A and $1.05$ under B. Under FIFO the wait is set by the second moment of the work,
not its mean, so less expected demand does not mean fewer jobs.

Under processor sharing a short job shares the server with the long one and leaves
after about $0.2$. With $\varphi\equiv1$ the mean number in system is $\rho/(1-\rho)$
for every work law (\cref{thm:insens} below): $1.0$ under A and $1.5$ under B, in the
order of the load. \Cref{thm:L-mono} together with insensitivity makes this true for
every capacity function $\varphi$.

\begin{figure}[t]
\centering
\begin{tikzpicture}[x=0.74cm, >=Stealth, font=\small, note/.style={font=\footnotesize, text=black!75}]
  \begin{scope}[yshift=0cm]
    \node[anchor=south west, font=\small, inner xsep=0pt] at (-0.9,2.79) {\textbf{Law A:} short jobs ($0.1$) with an occasional long one ($9.1$, probability $0.1$)};
    \fill[blue!15] (0,0.00) -- (0,0.00) -- (0,0.32) -- (1.5,0.32) -- (1.5,0.64) -- (2.2,0.64) -- (2.2,0.96) -- (4,0.96) -- (4,1.28) -- (5.5,1.28) -- (5.5,1.60) -- (6.3,1.60) -- (6.3,1.92) -- (8,1.92) -- (8,2.24) -- (9.1,2.24) -- (9.1,1.92) -- (9.2,1.92) -- (9.2,1.60) -- (9.3,1.60) -- (9.3,1.28) -- (9.4,1.28) -- (9.4,0.96) -- (9.5,0.96) -- (9.5,0.64) -- (9.6,0.64) -- (9.6,0.32) -- (9.6,0.32) -- (9.6,0.64) -- (9.7,0.64) -- (9.7,0.32) -- (9.8,0.32) -- (9.8,0.00) -- (10.2,0.00) -- (10.2,0.32) -- (10.3,0.32) -- (10.3,0.00) -- (11.8,0.00) -- (11.8,0.32) -- (11.9,0.32) -- (11.9,0.00) -- (13,0.00) -- (13,0) -- cycle;
    \draw[blue!70!black, thick] (0,0.00) -- (0,0.00) -- (0,0.32) -- (1.5,0.32) -- (1.5,0.64) -- (2.2,0.64) -- (2.2,0.96) -- (4,0.96) -- (4,1.28) -- (5.5,1.28) -- (5.5,1.60) -- (6.3,1.60) -- (6.3,1.92) -- (8,1.92) -- (8,2.24) -- (9.1,2.24) -- (9.1,1.92) -- (9.2,1.92) -- (9.2,1.60) -- (9.3,1.60) -- (9.3,1.28) -- (9.4,1.28) -- (9.4,0.96) -- (9.5,0.96) -- (9.5,0.64) -- (9.6,0.64) -- (9.6,0.32) -- (9.6,0.32) -- (9.6,0.64) -- (9.7,0.64) -- (9.7,0.32) -- (9.8,0.32) -- (9.8,0.00) -- (10.2,0.00) -- (10.2,0.32) -- (10.3,0.32) -- (10.3,0.00) -- (11.8,0.00) -- (11.8,0.32) -- (11.9,0.32) -- (11.9,0.00) -- (13,0.00);
    \draw[->] (0,0) -- (13.4,0) node[right, note] {$t$};
    \draw[->] (0,0) -- (0,2.54) node[left, note] {$N(t)$};
    \draw (-0.08,0.00) -- (0,0.00) node[left, note] {0};
    \draw (-0.08,0.64) -- (0,0.64) node[left, note] {2};
    \draw (-0.08,1.28) -- (0,1.28) node[left, note] {4};
    \draw (-0.08,1.92) -- (0,1.92) node[left, note] {6};
    \fill[red!55, draw=red!70!black, line width=0.3pt] (0,-0.55) rectangle (9.1,-0.25);
    \fill[blue!55, draw=blue!70!black, line width=0.3pt] (9.1,-0.55) rectangle (9.2,-0.25);
    \fill[blue!55, draw=blue!70!black, line width=0.3pt] (9.2,-0.55) rectangle (9.3,-0.25);
    \fill[blue!55, draw=blue!70!black, line width=0.3pt] (9.3,-0.55) rectangle (9.4,-0.25);
    \fill[blue!55, draw=blue!70!black, line width=0.3pt] (9.4,-0.55) rectangle (9.5,-0.25);
    \fill[blue!55, draw=blue!70!black, line width=0.3pt] (9.5,-0.55) rectangle (9.6,-0.25);
    \fill[blue!55, draw=blue!70!black, line width=0.3pt] (9.6,-0.55) rectangle (9.7,-0.25);
    \fill[blue!55, draw=blue!70!black, line width=0.3pt] (9.7,-0.55) rectangle (9.8,-0.25);
    \fill[blue!55, draw=blue!70!black, line width=0.3pt] (10.2,-0.55) rectangle (10.3,-0.25);
    \fill[blue!55, draw=blue!70!black, line width=0.3pt] (11.8,-0.55) rectangle (11.9,-0.25);
    \node[note, left] at (0,-0.4) {server};
    \fill (0,-0.62) -- ++(-0.09,-0.16) -- ++(0.18,0) -- cycle;
    \fill (1.5,-0.62) -- ++(-0.09,-0.16) -- ++(0.18,0) -- cycle;
    \fill (2.2,-0.62) -- ++(-0.09,-0.16) -- ++(0.18,0) -- cycle;
    \fill (4,-0.62) -- ++(-0.09,-0.16) -- ++(0.18,0) -- cycle;
    \fill (5.5,-0.62) -- ++(-0.09,-0.16) -- ++(0.18,0) -- cycle;
    \fill (6.3,-0.62) -- ++(-0.09,-0.16) -- ++(0.18,0) -- cycle;
    \fill (8,-0.62) -- ++(-0.09,-0.16) -- ++(0.18,0) -- cycle;
    \fill (9.6,-0.62) -- ++(-0.09,-0.16) -- ++(0.18,0) -- cycle;
    \fill (10.2,-0.62) -- ++(-0.09,-0.16) -- ++(0.18,0) -- cycle;
    \fill (11.8,-0.62) -- ++(-0.09,-0.16) -- ++(0.18,0) -- cycle;
    \node[anchor=west, draw=blue!50!black, fill=blue!5, rounded corners=1pt, font=\footnotesize, inner sep=3pt] at (13.9,1.12) {\begin{tabular}{@{}l@{\ \ }r@{}}work served & $10.0$ \\ mean $N(t)$ & $3.0$ \\ max $N(t)$ & $7$\end{tabular}};
  \end{scope}
  \begin{scope}[yshift=-3.35cm]
    \node[anchor=south west, font=\small, inner xsep=0pt] at (-0.9,1.19) {\textbf{Law B:} every job takes $1.2$ ($20\,\%$ more work per job)};
    \fill[orange!15] (0,0.00) -- (0,0.00) -- (0,0.32) -- (1.2,0.32) -- (1.2,0.00) -- (1.5,0.00) -- (1.5,0.32) -- (2.2,0.32) -- (2.2,0.64) -- (2.7,0.64) -- (2.7,0.32) -- (3.9,0.32) -- (3.9,0.00) -- (4,0.00) -- (4,0.32) -- (5.2,0.32) -- (5.2,0.00) -- (5.5,0.00) -- (5.5,0.32) -- (6.3,0.32) -- (6.3,0.64) -- (6.7,0.64) -- (6.7,0.32) -- (7.9,0.32) -- (7.9,0.00) -- (8,0.00) -- (8,0.32) -- (9.2,0.32) -- (9.2,0.00) -- (9.6,0.00) -- (9.6,0.32) -- (10.2,0.32) -- (10.2,0.64) -- (10.8,0.64) -- (10.8,0.32) -- (11.8,0.32) -- (11.8,0.64) -- (12,0.64) -- (12,0.32) -- (13,0.32) -- (13,0) -- cycle;
    \draw[orange!70!black, thick] (0,0.00) -- (0,0.00) -- (0,0.32) -- (1.2,0.32) -- (1.2,0.00) -- (1.5,0.00) -- (1.5,0.32) -- (2.2,0.32) -- (2.2,0.64) -- (2.7,0.64) -- (2.7,0.32) -- (3.9,0.32) -- (3.9,0.00) -- (4,0.00) -- (4,0.32) -- (5.2,0.32) -- (5.2,0.00) -- (5.5,0.00) -- (5.5,0.32) -- (6.3,0.32) -- (6.3,0.64) -- (6.7,0.64) -- (6.7,0.32) -- (7.9,0.32) -- (7.9,0.00) -- (8,0.00) -- (8,0.32) -- (9.2,0.32) -- (9.2,0.00) -- (9.6,0.00) -- (9.6,0.32) -- (10.2,0.32) -- (10.2,0.64) -- (10.8,0.64) -- (10.8,0.32) -- (11.8,0.32) -- (11.8,0.64) -- (12,0.64) -- (12,0.32) -- (13,0.32);
    \draw[->] (0,0) -- (13.4,0) node[right, note] {$t$};
    \draw[->] (0,0) -- (0,0.94) node[left, note] {$N(t)$};
    \draw (-0.08,0.00) -- (0,0.00) node[left, note] {0};
    \draw (-0.08,0.32) -- (0,0.32) node[left, note] {1};
    \draw (-0.08,0.64) -- (0,0.64) node[left, note] {2};
    \fill[orange!55, draw=orange!70!black, line width=0.3pt] (0,-0.55) rectangle (1.2,-0.25);
    \fill[orange!55, draw=orange!70!black, line width=0.3pt] (1.5,-0.55) rectangle (2.7,-0.25);
    \fill[orange!55, draw=orange!70!black, line width=0.3pt] (2.7,-0.55) rectangle (3.9,-0.25);
    \fill[orange!55, draw=orange!70!black, line width=0.3pt] (4,-0.55) rectangle (5.2,-0.25);
    \fill[orange!55, draw=orange!70!black, line width=0.3pt] (5.5,-0.55) rectangle (6.7,-0.25);
    \fill[orange!55, draw=orange!70!black, line width=0.3pt] (6.7,-0.55) rectangle (7.9,-0.25);
    \fill[orange!55, draw=orange!70!black, line width=0.3pt] (8,-0.55) rectangle (9.2,-0.25);
    \fill[orange!55, draw=orange!70!black, line width=0.3pt] (9.6,-0.55) rectangle (10.8,-0.25);
    \fill[orange!55, draw=orange!70!black, line width=0.3pt] (10.8,-0.55) rectangle (12,-0.25);
    \fill[orange!55, draw=orange!70!black, line width=0.3pt] (12,-0.55) rectangle (13,-0.25);
    \node[note, left] at (0,-0.4) {server};
    \fill (0,-0.62) -- ++(-0.09,-0.16) -- ++(0.18,0) -- cycle;
    \fill (1.5,-0.62) -- ++(-0.09,-0.16) -- ++(0.18,0) -- cycle;
    \fill (2.2,-0.62) -- ++(-0.09,-0.16) -- ++(0.18,0) -- cycle;
    \fill (4,-0.62) -- ++(-0.09,-0.16) -- ++(0.18,0) -- cycle;
    \fill (5.5,-0.62) -- ++(-0.09,-0.16) -- ++(0.18,0) -- cycle;
    \fill (6.3,-0.62) -- ++(-0.09,-0.16) -- ++(0.18,0) -- cycle;
    \fill (8,-0.62) -- ++(-0.09,-0.16) -- ++(0.18,0) -- cycle;
    \fill (9.6,-0.62) -- ++(-0.09,-0.16) -- ++(0.18,0) -- cycle;
    \fill (10.2,-0.62) -- ++(-0.09,-0.16) -- ++(0.18,0) -- cycle;
    \fill (11.8,-0.62) -- ++(-0.09,-0.16) -- ++(0.18,0) -- cycle;
    \node[anchor=west, draw=orange!50!black, fill=orange!5, rounded corners=1pt, font=\footnotesize, inner sep=3pt] at (13.9,0.32) {\begin{tabular}{@{}l@{\ \ }r@{}}work served & $12.0$ \\ mean $N(t)$ & $1.0$ \\ max $N(t)$ & $2$\end{tabular}};
  \end{scope}
  \node[note] at (6.5,-4.5) {$\blacktriangle$ arrival (the same ten arrival times in both panels)};
\end{tikzpicture}
\caption{One sample path of each work law at a FIFO server, on the same arrival times.
Top of each panel: the number in system $N(t)$; bottom: when the server works on each
job (red: the long job of law A). \emph{What to see:} law B does $20\,\%$ more work,
yet law A keeps about three times as many jobs waiting, because one long job blocks
everything behind it.}
\label{fig:more-work}
\end{figure}

\begin{theorem}[Monotonicity of the mean batch size]\label{thm:L-mono}
Let $\varphi$ be any positive function and $R$ the radius of convergence of the power
series $G(\rho)=\sum_n\rho^n/F(n)$. The mean $L(\rho)=\sum_n n\pi(n)$ of the stationary
distribution \eqref{L2:eq:piphi} is differentiable on $(0,R)$ and
\[
  \frac{\dd L}{\dd\rho}=\frac{\Var_\rho(N)}{\rho}\ \ge 0 .
\]
Hence $L$ is nondecreasing in the load. If $\varphi$ is nondecreasing, then
$R=\sup_n\varphi(n)$.
\end{theorem}
\begin{proof}
Inside its radius of convergence a power series can be differentiated term by term,
and the differentiated series has the same radius (a standard theorem of analysis).
With $a_n=1/F(n)$ we have $\rho G'(\rho)=\sum_n n a_n\rho^n$ and
$\rho^2G''(\rho)=\sum_n n(n-1)a_n\rho^n$, so
\[
  L=\frac{\rho G'}{G},\qquad \E_\rho[N(N-1)]=\frac{\rho^2G''}{G}.
\]
By the quotient rule,
\[
  \rho\,\frac{\dd L}{\dd\rho}
  =\frac{\rho G'}{G}+\frac{\rho^2G''}{G}-\Bigl(\frac{\rho G'}{G}\Bigr)^2
  =\E[N]+\E[N(N-1)]-\E[N]^2=\Var(N)\ge0 .
\]
If $\varphi$ is nondecreasing with $\varphi_\infty:=\sup\varphi$, then
$a_{n+1}/a_n=1/\varphi(n+1)\to1/\varphi_\infty$, so by the ratio test
$R=\varphi_\infty$ ($R=\infty$ if $\varphi_\infty=\infty$).
\end{proof}

For $\rho>0$, $N$ is not degenerate, so in fact $L$ is strictly increasing. The theorem
says, without knowing anything about the shape of $\varphi$, that ``a policy that
increases the mean demand does not shrink the batch''. Since $\varphi$ must be
calibrated from measurements, such shape-free conclusions matter.

\section{Insensitivity}

\Cref{prop:phi} assumed exponential work. But the decode demand is proportional to the
output length, and there is no reason for it to be exponential. Remarkably, under PS
the conclusion does not depend on the work distribution.

\begin{definition}[Insensitivity]\label{L2:def:insens}
A stationary quantity of a queue (here, the distribution of the number of jobs) is
\term{insensitive} to the work distribution $G$ if it depends on $G$ only through its
mean $\int x\dd G(x)$.
\end{definition}

\begin{theorem}[Insensitivity of PS \cite{kelly1979,bcmp1975}]\label{thm:insens}
For a PS server with capacity $\varphi$ and Poisson arrivals of rate $\lambda$, the
stationary distribution of the number of jobs is insensitive to the work distribution
and equals \eqref{L2:eq:piphi}. We prove it here when the work is a finite mixture of
exponentials, that is, $\mathrm{Exp}(\mu_k)$ with probability $\alpha_k$
($k=1,\dots,K$, $\sum_k\alpha_k=1$).
\end{theorem}
\begin{proof}
Think of each job as independently drawing a type $k$ with probability $\alpha_k$ on
arrival and then drawing its work from $\mathrm{Exp}(\mu_k)$. The vector of job counts
by type, $\mathbf n=(n_1,\dots,n_K)$ with $n=|\mathbf n|=\sum_kn_k$, is a CTMC. By
thinning of the Poisson process (\cref{prop:superpose}(ii)), type $k$ arrives at rate
$\lambda\alpha_k$ ($\mathbf n\to\mathbf n+e_k$), and as in the proof of \cref{prop:phi}
a single type-$k$ job finishes at rate $\mu_k\varphi(n)/n$, so the rate of
$\mathbf n\to\mathbf n-e_k$ is $n_k\mu_k\varphi(n)/n$.
Let $\rho_k=\lambda\alpha_k/\mu_k$. We show that
\[
  \pi(\mathbf n)=B\,\frac{n!}{n_1!\cdots n_K!}\,\frac{\prod_k\rho_k^{n_k}}{F(n)}
\]
satisfies global balance. If $n_k\ge1$, then
$\pi(\mathbf n)/\pi(\mathbf n-e_k)=\dfrac{n}{n_k}\cdot\dfrac{\rho_k}{\varphi(n)}$, so
\begin{equation}\label{L2:eq:partial}
  \pi(\mathbf n)\,\frac{n_k\mu_k\varphi(n)}{n}=\pi(\mathbf n-e_k)\,\rho_k\mu_k=\pi(\mathbf n-e_k)\,\lambda\alpha_k .
\end{equation}
This says ``the flow out of $\mathbf n$ by a type-$k$ completion equals the flow into
$\mathbf n$ by a type-$k$ arrival''. Applying \eqref{L2:eq:partial} at
$\mathbf n+e_k$ and summing over $k$,
\[
  \sum_k\pi(\mathbf n+e_k)\,\frac{(n_k+1)\mu_k\varphi(n+1)}{n+1}=\sum_k\pi(\mathbf n)\lambda\alpha_k=\pi(\mathbf n)\lambda ,
\]
that is, ``the flow out by arrivals equals the flow in by completions''. Adding the two
gives global balance at $\mathbf n$. Splitting global balance in this way is called
\term{partial balance}. Now, by the multinomial theorem,
\[
  \Prob(N=n)=\sum_{|\mathbf n|=n}\pi(\mathbf n)=\frac{B}{F(n)}\sum_{|\mathbf n|=n}\frac{n!}{\prod_kn_k!}\prod_k\rho_k^{n_k}
  =\frac{B\,\rho^n}{F(n)},\qquad \rho=\sum_k\rho_k=\lambda\sum_k\frac{\alpha_k}{\mu_k}=\lambda\E[S].
\]
The normalisability condition is $G(\rho)<\infty$, with $B=1/G(\rho)$. This is
\eqref{L2:eq:piphi}, and $G$ involves only $\E[S]$.
(Uniqueness follows from irreducibility and \cref{L2:thm:ergodic}.)
\end{proof}

\begin{remark}
Insensitivity for general work distributions is due to \cite{kelly1979} and
\cite{bcmp1975} (BCMP). One route to a proof is to extend the partial balance argument
above to phase-type distributions and take limits, using the fact that phase-type
distributions are dense, under weak convergence, among distributions on $[0,\infty)$.
Here we only cite it.
The same insensitivity holds for a \term{delay station} (an infinite-server station
where every job immediately gets its own server, so nobody waits), which we use in
\cref{sec:L4}.
\end{remark}

\paragraph{FIFO is not insensitive.} Compare three service-time distributions with the
same mean $\E[S]=1$ and the same load $\rho=0.5$ ($\lambda=0.5$). Borrowing from
\cref{sec:L3} the Pollaczek--Khinchine formula $W_q=\lambda\E[S^2]/(2(1-\rho))$ and
Little's law $L=\lambda W_q+\rho$, we get the following.
\begin{center}\small
\begin{tabular}{@{}lccc@{}}
\toprule
Service time $S$ & $\E[S^2]$ & FIFO $L$ & PS $L$ \\
\midrule
constant 1 & 1 & 0.75 & 1 \\
$\mathrm{Exp}(1)$ & 2 & 1.00 & 1 \\
$0.1$ w.p.\ $0.9$, $9.1$ w.p.\ $0.1$ & 8.29 & 2.57 & 1 \\
\bottomrule
\end{tabular}
\end{center}
Under FIFO one long job blocks every job behind it (head-of-line blocking). Under PS
short jobs share capacity with the long one and leave first.
The fairness of PS is also visible job by job: in M/G/1-PS, the mean sojourn time of
a job with work $x$ is $x/(1-\rho)$, proportional to its own size and independent of
the size distribution of the other jobs \cite{kleinrock1976}.
This difference separates the two stages of the paper: the FIFO prefill queue sees the
second moment of the work, while the PS decode stage sees only the mean.

\section{The price at the decode stage}

\begin{proposition}[The decode stage prices only work]\label{prop:decode-price}
For a PS station with constant capacity $C$ and load $\rho<C$, $L=\rho/(C-\rho)$. If a
policy change increases the work of a fraction $q_i$ of arrivals by $\Delta D_i\ge0$,
raising the load to $\rho'=\rho+\lambda\sum_iq_i\Delta D_i<C$, then exactly
\[
  \Delta L=\frac{C(\rho'-\rho)}{(C-\rho)(C-\rho')}
  =\frac{C-\rho}{C-\rho'}\;\lambda\sum_iq_i\,\Phi^D_i,\qquad
  \Phi^D_i=\frac{C\,\Delta D_i}{(C-\rho)^2} .
\]
In particular $\lambda\sum_iq_i\Phi^D_i\le\Delta L$.
\end{proposition}
\begin{proof}
With $\varphi\equiv C$ we have $F(n)=C^n$, and \eqref{L2:eq:piphi} is geometric with
ratio $\rho/C$, so $L=(\rho/C)/(1-\rho/C)=\rho/(C-\rho)$. Therefore
\[
  \Delta L=\frac{\rho'}{C-\rho'}-\frac{\rho}{C-\rho}=\frac{\rho'(C-\rho)-\rho(C-\rho')}{(C-\rho)(C-\rho')}
  =\frac{C(\rho'-\rho)}{(C-\rho)(C-\rho')} .
\]
Substituting $\rho'-\rho=\lambda\sum_iq_i\Delta D_i$ and multiplying and dividing by
$(C-\rho)^2$ gives the second equality. Since $\rho'\ge\rho$, the leading factor is at
least one.
\end{proof}

$\Phi^D_i=\Delta D_i\,L'(\rho)$ is the first-order approximation, and the factor
$(C-\rho)/(C-\rho')$ corrects for curvature. For example, with $C=1$ and
$\rho=0.5\to0.6$, $\Delta L=0.5$, the first-order term is $0.4$, and the factor is
$1.25$.

\begin{corollary}\label{L2:cor:decode-hit}
If a policy change alters the distribution of the decode work but not its mean $\E[D]$,
the distribution of the number of jobs at the decode stage does not change.
\end{corollary}
\begin{proof}
By \cref{thm:insens} the distribution depends only on $\rho_D=\lambda\E[D]$.
\end{proof}

\serving{KV-cache hits and misses change only the prefill work; they do not change the
decode demand $D_i=o_i(\beta K_i+\omega/m)$. So by \cref{L2:cor:decode-hit} a miss costs
nothing at the decode stage to first order, and its whole price is paid at the FIFO
prefill queue (\cref{sec:L3}). Only decisions that add decode demand, such as admitting
a session, pay the load factor of \cref{prop:decode-price}.
\Cref{prop:decode-price,L2:cor:decode-hit} are the content of the paper's proposition
``The decode stage prices only work''.}

\section{Exercises}

\begin{exercise}\label{L2:ex:mminf}
For $\varphi(n)=n$ (each job has its own server, M/M/$\infty$), show that
\eqref{L2:eq:piphi} is the Poisson distribution with mean $\rho$, and check the formula
$\dd L/\dd\rho=\Var(N)/\rho$ of \cref{thm:L-mono} directly.
\emph{Hint:} $F(n)=n!$, and a Poisson distribution has equal mean and variance.
\end{exercise}

\begin{exercise}\label{L2:ex:mm1wait}
Show that in M/M/1-FIFO the sojourn time of an arriving job is $\mathrm{Exp}(\mu-\lambda)$.
\emph{Hint:} by PASTA (\cref{thm:pasta}) the number found on arrival is distributed as
$(1-\rho)\rho^n$, and conditionally the sojourn time is a sum of $n+1$ independent
$\mathrm{Exp}(\mu)$ variables.
\end{exercise}

\begin{exercise}[M/M/$m$ and the Erlang C formula]\label{L2:ex:mmm}
Arrivals are Poisson with rate $\lambda$, service times are $\mathrm{Exp}(\mu)$, and there
are $m$ servers; a customer waits only when all $m$ are busy. With $a=\lambda/\mu<m$,
show from \cref{thm:birthdeath} that
$\pi(n)=\pi(0)\,a^n/n!$ for $n\le m$ and $\pi(n)=\pi(0)\,a^n/(m!\,m^{n-m})$ for $n\ge m$,
that the probability of waiting is
\[
  C(m,a)=\frac{\frac{a^m}{m!}\frac{m}{m-a}}{\sum_{k=0}^{m-1}\frac{a^k}{k!}+\frac{a^m}{m!}\frac{m}{m-a}},
\]
and that the mean wait is $W_q=C(m,a)/(m\mu-\lambda)$. This is the memory-slot station
of \cref{L1:sec:projection} when every turn holds one slot for an exponential time.
\emph{Hint:} the death rate in state $n$ is $\min(n,m)\,\mu$; use PASTA for the waiting
probability and Little's law for $W_q$.
\end{exercise}
